\documentclass[a4paper,12pt]{article}
\usepackage[utf8]{inputenc}
\usepackage{geometry}
\usepackage{graphicx}   % images
\usepackage{fancyhdr}   % headers/footers
\usepackage{tcolorbox}
\usepackage{listings}
\usepackage{xcolor}
\geometry{margin=1in}

% ---------- Header ----------
\setlength{\headheight}{36pt}
\setlength{\headsep}{18pt}
\renewcommand{\headrulewidth}{0.4pt}
\fancyhf{}
\fancyhead[L]{\includegraphics[width=0.13\textwidth, keepaspectratio]{Figures/UM6Plogo.png}}
\fancyhead[R]{\includegraphics[width=0.13\textwidth, keepaspectratio]{Figures/CC.jpg}}
\fancyfoot[L]{Data Management Lab}
\fancyfoot[R]{Prof. Karima Echihabi}
\fancyfoot[C]{Page \thepage}

% ---------- Deliverable Template ----------
\begin{document}
\thispagestyle{empty}
\begin{center}
  \includegraphics[width=0.25\textwidth]{Figures/UM6Plogo.png}\hfill
  \includegraphics[width=0.25\textwidth]{Figures/CC.jpg}
  \vspace{1.2cm}

  {\LARGE \textbf{Deliverable : \:Logical design Lab2 }}\\[0.6cm]
  {\large \textbf{Data Management Course}}\\[0.2cm]
  {\large UM6P College of Computing}\\[0.8cm]

  {\normalsize \textbf{Professor:} Karima Echihabi \quad 
   \textbf{Program:} Computer Engineering}\\[0.1cm]
  {\normalsize \textbf{Session:} Fall 2026}\\[1cm]

  \rule{0.9\textwidth}{0.5pt}\\[0.5cm]
  {\large \textbf{Team Information}} \\[0.3cm]
  \begin{tabular}{|l|l|}
    \hline
    \textbf{Team Name} & Team N°1 \\ \hline
    \textbf{Member 1}  & Nizar El HAZZITI  \\ \hline
    \textbf{Member 2}  & Mohamed JBARI  \\ \hline
    \textbf{Member 3}  & Riyad HABIBI \\ \hline
    \textbf{Member 4}  & Chaimae IAA    \\ \hline
    \textbf{Repository Link} & \texttt{https://github.com/...} \\ \hline
  \end{tabular}
  \rule{0.9\textwidth}{0.5pt}\\
\end{center}
\clearpage
\pagestyle{fancy}

% ---------- Sections for Students ----------
\section{Introduction}
This lab focuses on the logical design phase of the RLMS project,which consists of translating the ER diagram built in Lab 1 into actual database tables with rows, columns and primary keys while preserving the semantics captured in the conceptual model.

\section{Requirements}
This deliverable translates the RLMS ER diagram into a relational schema, mapping all entity sets, relationships (including weak entities, ISA hierarchies, recursive and ternary relationships, and aggregation) into properly keyed tables ,specifying primary keys, foreign keys. It also identifies any constraints not expressible via keys alone, and delivers the full .sql schema plus a short report on  the team, challenges, and design choices.

\section{Team roles}
Nizar: Modelizing relational table for entities Person and Laboratories.

\section{Key challenges}

-For Key constraints we won't need and individual table for such a relationship, therefore we add one an attribute to the entitie, otherwhise we create another table for the relationalship.\\
-For M:1 relationships (like monitors) no need to create a new table , we just add an attribute technicalId to the laboStock relationship , since each stock has only one monitor.\\
-"Supplier" is a  ternary relationship with attribute "unitPrice", the attribute then depends on the combination of the three entities (lab, supplier, consumable) , so we can't split it into labSupp and consumSupp.\\


\section{Discussion}
Lessons learned:\\
- We don't split a ternary relationship into two tables ! It stays one table with the 3 Foreign Keys .\\
-When merging a relationship into another table , we NEVER change the primary key (like monitors).\\
-A relationship based on top of an aggregated relationship (like stocks) must reference that aggregate's one table NOT the original entities it was built from.\\


\section{Conclusion and toughts for next lab}
A simple .sql file where we implementated the relational model. The only thing to address better is the consistency of code style, because the each member has his style (example: the use of camel or snake case, tabulations ...), we should find a good style and all of us use it.


\end{document}
